%%%%%%%%%%%%%%%%%%%%%%%%%%%%%%%%%%%%%%%%%%%%%%%%%%%%%%%%%%%%%%%%%%%%%%%%%%%%%%%%%
\documentclass{article}
% Feel free to contact me for any reason:
%% Email 1: zain.kamal@rutgers.edu
%% Email 2: z.kamal2021@gmail.com
%% Discord: alci#6038

% Last updated 1/28/22

% Template based off of Justin Kim's (don't know where he got it from originally), but I've made a ton of edits. His soul lives on in random packages and commands I'm too lazy to comment out.


%%%%%%%%%%%%%%%%%%%%%%%%%%%%%%%%%%%%%%%%%%%%%%%%%%%%%%%%%%%%%%%%%%%%%%%%%%%%%%%%%%%%
%%% Default packages:

\usepackage[margin=1in]{geometry} 
\usepackage{amsmath,amsthm,amssymb,amsfonts, fancyhdr, color, comment, graphicx, environ}
\usepackage{xcolor}
\usepackage{mdframed}
\usepackage{bm}
\usepackage[shortlabels]{enumitem}
\usepackage{mathtools}
\usepackage{listings}
\usepackage{stmaryrd}
\usepackage{indentfirst}
\usepackage{hyperref}


%%%%%%%%%%%%%%%%%%%%%%%%%%%%%%%%%%%%%%%%%
%%% Pacakages I've added:

%% For Math 300 (Winter 2021):

\usepackage{mathdots} % for \iddots

\usepackage[ruled]{algorithm2e} % Algorithms
% NOTE: FIND A BETTER ALGORITHM PACKAGE, OR ATLEAST LEARN HOW TO USE THIS ONE BECAUSE DOUBLE INDENTING IS A FUCKING NIGHTMARE (or just import from mathcha?)
%% Example algorithm:
% \begin{center}
% 	\begin{minipage}{0.5\linewidth} % Adjust the minipage width to accomodate for the length of algorithm lines
% 		\begin{algorithm}[H]
% 			\KwIn{$(a, b)$, two floating-point numbers}  % Algorithm inputs
% 			\KwResult{$(c, d)$, such that $a+b = c + d$} % Algorithm outputs/results
% 			\medskip
% 			\If{$\vert b\vert > \vert a\vert$}{
% 				exchange $a$ and $b$ \;
% 			}
% 			$c \leftarrow a + b$ \;
% 			$z \leftarrow c - a$ \;
% 			$d \leftarrow b - z$ \;
% 			{\bf return} $(c,d)$ \;
% 			\caption{\texttt{FastTwoSum}} % Algorithm name
% 			\label{alg:fastTwoSum}   % optional label to refer to
% 		\end{algorithm}
% 	\end{minipage}
% \end{center}



%%%%%%%%%%%%%%%%%%%%%%%%%%%%%%%%%%%%%%%%%%%%%%%%%%%%%%%%%%%%%%%%%%%%%%%%%%%%%%%%%%%%
%%% Default commands:

\renewcommand{\vec}[1]{\mathbf{#1}}
	
\newcommand{\WidestEntry}{$lon_1$}%
\newcommand{\SetToWidest}[1]{\makebox[\widthof{\WidestEntry}]{#1}}%
\newcommand\tab[1][0.61cm]{\hspace*{#1}}
\newcommand{\nats}{\mathbb{N}}
\newcommand{\rats}{\mathbb{Q}}
\newcommand{\reals}{\mathbb{R}}
\newcommand{\Z}[1]{\mathbb{Z}_{#1}}
\newcommand{\BigO}[1]{\mathcal{O}(#1)}
\newcommand{\seq[1]}{(#1_n)}
\newcommand{\subseq[1]}{(#1_{n_k})}
\newcommand{\Lim}[2]{\lim \limits _{#1 \to #2}}
\newcommand{\Min}[2]{\min \{#1, #2\}}
\newcommand{\inv}{^{-1}}
\newcommand{\h}{^\text{th}}
\newcommand{\lrangle}[1]{\langle #1 \rangle}
\newcommand{\abs}[1]{\left\lvert #1 \right\rvert}

\DeclarePairedDelimiter{\ceil}{\lceil}{\rceil}
\DeclarePairedDelimiter{\floor}{\lfloor}{\rfloor}
\DeclareMathOperator{\supp}{supp}
\DeclareMathOperator{\rad}{rad}
\DeclareMathOperator*{\argmin}{arg\,min}
\DeclareMathOperator*{\argmax}{arg\,max}
\DeclareMathOperator*{\Var}{Var}
\DeclareMathOperator*{\Cov}{Cov}
\DeclareMathOperator*{\Corr}{Corr}
\DeclareMathOperator*{\Aut}{Aut}
\newcommand{\prob}[1]{\section*{Problem #1}}


%%%%%%%%%%%%%%%%%%%%%%%%%%%%%%%%%%%%%%%%%
%%% Commands I've added:

%% For Math 300 (Winter 2021):

\newcommand{\lrbrace}[1]{\{ #1 \}}
\newcommand{\powerset}{\mathcal{P}}
\newcommand{\ints}{\mathbb{Z}}

% Source/inspiration: https://tex.stackexchange.com/a/42728:
\newcommand{\numberthis}{\addtocounter{equation}{1}\tag{\theequation}\label{\theequation}}
    % Within an `align*` environment, put `\numberthis` after a line to number it. 
    % Access it with `\eqref{ [number of equation] }`
\newcommand{\numberthiswith}[1]{\addtocounter{equation}{1}\tag{\theequation}\label{#1}}
    % Within an `align*` environment, put `\numberthiswith{ [your_label] }` after a line to number it. 
    % Access it with `\eqref{ [your_label] }`



%%%%%%%%%%%%%%%%%%%%%%%%%%%%%%%%%%%%%%%%%%%%%%%%%%%%%%%%%%%%%%%%%%%%%%%%%%%%%%%%%%%%
%%% Default formatting:

\hypersetup{
    colorlinks=true,
    linkcolor=blue,
    filecolor=magenta,      
    urlcolor=blue,
}

\setlength{\parindent}{0cm}
\setlength{\parskip}{6pt}

\pagestyle{fancy}


%% Misc formatting additions

% make bullets with itemize much smaller
\newlength{\mylen}
\setbox1=\hbox{$\bullet$}\setbox2=\hbox{\tiny$\bullet$}
\setlength{\mylen}{\dimexpr0.5\ht1-0.5\ht2}
\renewcommand\labelitemi{\raisebox{\mylen}{\tiny$\bullet$}}


% Modified version of problem environment below
% \newenvironment{problem}[2][Problem]
%     { \begin{mdframed}[backgroundcolor=gray!5] \textbf{#1 #2} \\}
%     {  \end{mdframed}}
% \newenvironment{solution}{\textbf{Solution}\\}


%%%%%%%%%%%%%%%%%%%%%%%%%%%%%%%%%%%%%%%%%
%%% Formatting I've added:

%% Grey boxes for problem statements (note that I don't have a "solution" section):

% Problem environment, but shows "(a)" instead of "Problem a"
\newenvironment{problem}[2][]
    { \begin{mdframed}[backgroundcolor=gray!5] \textbf{#1 (#2)}}
    {  \end{mdframed}}
% Problem environment, but no "([input_char])" at all
\newenvironment{problem*}
    { \begin{mdframed}[backgroundcolor=gray!5] \\}
    {  \end{mdframed}}


% Example environment, currently identical to "problem" (Note: this is better written than the problem environment because I wrote it myself from scratch. Use this as an example for future new environments.)
\newcounter{example}[section]
\newenvironment{example}
    { 
        \refstepcounter{example}
        \begin{mdframed}[backgroundcolor=gray!5]
        \textbf{\\Example \thesection.\theexample:}
    }
    {\\ \end{mdframed}}


%%%%%%%%%%%%%%%%%%%%%%%%%%%%%%%%%%%%%%%%%%%%%%%%%%%%%%%%%%%%%%%%%%%%%%%%%%%%%%%%%%%%
% Misc things I've added




%%%%%%%%%%%%%%%%%%%%%%%%%%%%%%%%%%%%%%%%%%%%%
% Fill in the appropriate information below
\lhead{Zain Kamal}
% \rhead{Math 244 Spring 2022} % Moved to document
% \chead{\textbf{Homework 2}} % Moved to document

\usepackage{tikz}
\usepackage{tikz-among-us}



\begin{document}
\chead{\textbf{Homework 4}}
%%%%%%%%%%%%%%%%%%%%%%%%%%%%%%%%%%%%%%%%%%%%%%%%%%%%%%%%%%%%%%%%%%%%%%%%%%%%%%%%%
\section{}

\begin{problem*}
    Using induction, prove that \[\sum_{i=1}^n i^2=\frac{n(n+1)(2n+1)}{6}.\]
\end{problem*}

Define $F(n):= \left( \sum_{i=1}^n i^2=\frac{n(n+1)(2n+1)}{6} \right)$. We aim to prove by induction that $F(n)$ is true for all $n \in \nats$.

\

For the base case, we evaluate $F(1)$ is true because $1^2=\tfrac{(1)(2)(3)}{6}$.

\

% For the induction step, we assume $F(n)$ is true for $n=1$ and prove $F(n+1)$ is true for all $n \in \nats$. 
For the induction step, we aim to prove that if we assume $F(k)$ is true for an arbitrary natural number $k \geq 1$, then $F(k+1)$ is also true. 
% For the induction step, we assume $F(n)$ is true for all $n \in \nats$ and prove $F(n+1)$ is true. 
We begin by expanding the left side (arithmetic sequence) of $F(k+1)$:
\begin{align*}
    \sum_{i=1}^{k+1} i^{2} &=(k+1)^{2}+\left(\sum_{i=1}^{k} i^{2}\right) \\
    &=\left(k^{2}+2 k+1\right)+\tfrac{1}{6}(k(k+1)(2 k+1)) \quad \quad (\textrm{Using the definition of } F(k)) \\
    &=\tfrac{1}{6}\left(6 k^{2}+12 k+6\right)+\tfrac{1}{6}\left(2 k^{3}+3 k^{2}+k\right) \\
    &=\tfrac{1}{6}\left(2 k^{3}+9 k^{2}+13 k+6\right). \numberthis
\end{align*}
Next, we expand the right side (fraction) of $F(k+1)$:
\begin{align*} 
    & \frac{(k+1)(k+1+1)(2(k+1)+1)}{6} \\
    =& \frac{1}{6}\left(2 k^{3}+9 k^{2}+13 k+6\right). \numberthis
\end{align*}
We see that equation (1) equals equation (2).
% , so we can conclude that if we assume $F(k)$ is true for an arbitrary natural number $k \geq 1$, then $F(k+1)$ is also true.

\

Having proven $F(1)$ is true, and that if we assume $F(k)$ is true for an arbitrary natural number $k \geq 1$ then $F(k+1)$ is also true, we can use the principle of induction to conclude $F(n)$ is true for all $n \in \nats$. \qed

\
\hline
%%%%%%%%%%%%%%%%%%%%%%%%%%%%%%%%%%%%%%%%%%%%%%%%%%%%%%%%%%%%%%%%%%%%%%%%%%%%%%%%%
\section{}

\begin{problem*}
    Using induction, show that for all $n\in\nats,$ 8 divides $5^{2n}-1.$
\end{problem*}

Using the definition of ``divides'' from Definition 3.19 in Lecture Notes 3, we represent ``8 divides $5^{2n}-1$'' by defining the proposition
\begin{equation*}
    F(n):= \textrm{there is some integer } c \textrm{ such that }  8c=5^{2n}-1.
\end{equation*}
We aim to prove by induction that $F(n)$ is true for all $n \in \nats$.

\

For the base case, we evaluate $F(1)$ by solving for some $c_1$ such that $8c_1&=5^{2(1)}-1$. We find that $c_1=3$, which is an integer, so $F(1)$ is true.

\

For the induction step, we aim to prove that if we assume $F(k)$ is true for an arbitrary natural number $k \geq 1$, then $F(k+1)$ is also true. 

The initial assumption that $F(k)$ is true implies that there exists some integer $c_2$ such that $8c_2=5^{2(k)}-1$. Some basic manipulation yields $5^{2k}=8c_2+1$, which will be useful in the next step.

We proceed to evalute $F(k+1)$ by solving for some $c_3$ such that $8c_3 = 5^{2(k+1)}-1$. Some basic manipuation yields:
\begin{align*}
    8c_3 &= 5^{2(k+1)}-1 \\
    &= 5^{2k+2}-1 \\
    &= (5^{2k})(25)-1,
\end{align*}
at which point we can substitue $5^{2k}=8c_2+1$ to find:
\begin{align*}
    8c_3 &= (8c_2+1)(25)-1 \\
    &= 200c_2+24\\
    &= 8(25c_2+3)\\
    c_3 &= 25c_2+3.
\end{align*}

Because $c_2$ is an integer, it follows that $c_3$ is also an integer. Therefore if we assume $F(k)$ is true for an arbitrary natural number $k \geq 1$, then $F(k+1)$ is also true. 

\

Having proven $F(1)$ is true, and that if we assume $F(k)$ is true for an arbitrary natural number $k \geq 1$ then $F(k+1)$ is also true, we can use the principle of induction to conclude $F(n)$ is true for all $n \in \nats$. \qed

\
\hline
%%%%%%%%%%%%%%%%%%%%%%%%%%%%%%%%%%%%%%%%%%%%%%%%%%%%%%%%%%%%%%%%%%%%%%%%%%%%%%%%%
\section{}

\begin{problem*}
    Suppose there are $n$ people in a room. Everyone in the room is very friendly, as it happens, and each person wants to shake every other person's hand. Show that there are $\frac{n^2-n}{2}$ handshakes that will occur. (For example, for $n=2$, there will be one handshake that occurs between the two people). 
\end{problem*}

Define the proposition 
\begin{equation*}
F( n) :=\ \text{for} \ n\ \text{people in a room,} \ \frac{n^{2} -n}{2} \ \text{handshakes will occur} .
\end{equation*}
We aim to prove by induction that $F(n)$ is true for all $n \in \nats$.

\

For the base case ($n=1$), we have 1 person in a room who shakes 0 hands. We also see that $\frac{( 1)^{2} -( 1)}{2} =0$, so $F( 1)$ is true.

\

For the induction step, we aim to prove that if we assume $F( k)$ is true for an arbitrary natural number $k\geq 1$, then $F( k+1)$ is also true.

The assumption that $F(k)$ is true means that in a room with $k$ people, $\frac{k^{2} -k}{2}$ handshakes will occur. However, if one additional person enters the room, that new person will need to shake $k$ hands for every person already in the room. Thus, the number of handshakes increases to:
\begin{align*}
\frac{k^{2} -k}{2} +k & =\frac{k^{2} -k+2k}{2}\\
 & =\frac{k^{2} +k}{2} . \numberthis
\end{align*}
We also see that substituting $n=k+1$ into the formula given by the question yields:
\begin{align*}
\frac{( k+1)^{2} -( k+1)}{2} & =\frac{k^{2} +2k+1-k-1}{2}\\
 & =\frac{k^{2} +k}{2} ,
\end{align*}
which matches the number of handshakes in equation \eqref{3}. Therefore, if we assume $F( k)$ is true for an arbitrary natural number $k\geq 1$, then $F( k+1)$ is also true.

\

Having proven that $F( 1)$ is true, and that if we assume $F( k)$ is true for an arbitrary natural number $k\geq 1$ then $F( k+1)$ is also true, we can use the principle of induction to conclude that $F( n)$ is true for all $n\in \mathbb{N}$. \qed

\
\hline
%%%%%%%%%%%%%%%%%%%%%%%%%%%%%%%%%%%%%%%%%%%%%%%%%%%%%%%%%%%%%%%%%%%%%%%%%%%%%%%%%
\section{}

\begin{problem*}
    \textbf{(Towers of Hanoi).} Consider the following figure, where we have three rods, and 5 disks stacked on top of each other on the leftmost rod, and the largest disk is on the bottom, and the smallest is on top. The objective of the puzzle is to move all the disks to the rightmost rod, while observing the following three rules: \begin{itemize}
        \item Only one disk may be moved at a time. 
        \item A ``move" is taking the topmost disk from one of the rods and moving it to a different rod.
        \item A larger disk may not be placed on a smaller disk. 
    \end{itemize}
    Now consider the same game with $n$ disks instead of 5. Show, using induction, that the puzzle may be solved in $2^n-1$ moves.
\end{problem*}

Define the proposition 
\begin{equation*}
F( n) :=\ \text{a Towers of Hanoi puzzle with} \ n\ \text{disks may be solved in} \ 2^{n} -1\ \text{moves} .
\end{equation*}
We aim to prove by induction that $F(n)$ is true for all $n \in \nats$.

\

For the base case of one disk ($n=1$), it will take one move from the leftmost rod to the rightmost rod to solve the puzzle. We also see that $2^{( 1)} -1=1$, so $F( 1)$ is true.

\

For the induction step, we aim to prove that if we assume $F( k)$ is true for an arbitrary natural number $k\geq 1$, then $F( k+1)$ is also true.

In a scenario with $k+1$ disks, we can split the solution into three steps:
\begin{enumerate}
\item Since we assume $F( k)$ is true, it will take $2^{k} -1$ moves to move the $k$ smallest (top-most) disks from the leftmost rod to the middle rod. We are essentially solving the puzzle, but swapping the middle rod for the rightmost rod.
\item It will take one move to move the remaining disk on the leftmost rod (the largest disk) to the rightmost rod.
\item It will take another $2^{k} -1$ moves to move the $k$ disks on the middle rod to the rightmost rod. We are essentially reversing the first step, but swapping the leftmost rod for the rightmost rod.
\end{enumerate}

After these three steps, the total number of moves would be:
\begin{align*}
\left( 2^{k} -1\right) +( 1) +\left( 2^{k} -1\right) & =\left( 2^{k}\right)( 2) -1\\
 & =2^{k+1} -1.
\end{align*}
Notice that this is equal to $2^{n} -1$ for $n=k+1$. Therefore, if we assume $F(k)$ is true for an arbitrary natural number $k\geq 1$, then $F(k+1)$ is also true.

\

Having proven that $F(1)$ is true, and that if we assume $F(k)$ is true for an arbitrary natural number $k\geq 1$ then $F(k+1)$ is also true, we can use the principle of induction to conclude that $F(n)$ is true for all $n\in \mathbb{N}$.\qed 

\
\hline
%%%%%%%%%%%%%%%%%%%%%%%%%%%%%%%%%%%%%%%%%%%%%%%%%%%%%%%%%%%%%%%%%%%%%%%%%%%%%%%%%
\section{}

\begin{problem*}
    Show that for any convex polygon with $n$ vertices, the sum of interior angles is $(n-2)\pi$. A \textbf{convex polygon} is a polygon where for any two points in the interior, the line segment that joins them is contained within the polygon. However, you may use the easier definition, which is that a polygon is convex when all interior angles are less than $\pi$.
\end{problem*}

Define the proposition
\begin{equation*}
F( n) :=\text{in a convex polygon with} \ n\ \text{verticies, the sum of the interior angles is} \ \pi ( n-2) .
\end{equation*}
Because the minimum number of verticies in a polygon is $3$, we aim to prove by induction that $F(n)$ is true for all natural numbers $n \geq 3$.

\

For the base case of a triangle ($n=3$), consider the diagram below. Let the angles of the triangle be denoted by $x$, $y$, and $z$.
\begin{center}
    \tikzset{every picture/.style={line width=0.75pt}} %set default line width to 0.75pt
    \begin{tikzpicture}[x=0.75pt,y=0.75pt,yscale=-1,xscale=1]
        %uncomment if require: \path (0,300); %set diagram left start at 0, and has height of 300
        
        %Straight Lines [id:da6627804214979913] 
        \draw    (0,1) -- (280,1) ;
        %Shape: Triangle [id:dp5598900021195212] 
        \draw   (135,1) -- (220,91) -- (50,91) -- cycle ;
        
        \draw   (129.4,86) -- (139.4,91) -- (129.4,96) ;
        \draw   (84.6,-4) -- (94.6,1) -- (84.6,6) ;
        \draw   (180.6,-3.6) -- (190.6,1.4) -- (180.6,6.4) ;    
        
        % Text Node
        \draw (73,75) node [anchor=north west][inner sep=0.75pt]   [align=left] {x};
        % Text Node
        \draw (186,75) node [anchor=north west][inner sep=0.75pt]   [align=left] {y};
        % Text Node
        \draw (130,12) node [anchor=north west][inner sep=0.75pt]   [align=left] {z};
        % Text Node
        \draw (101,6) node [anchor=north west][inner sep=0.75pt]   [align=left] {x'};
        % Text Node
        \draw (158,6) node [anchor=north west][inner sep=0.75pt]   [align=left] {y'};
    
    \end{tikzpicture}
\end{center}
Angles $x'$, $y'$, and $z$ form a straight line with an angle measure of $\pi $. The top-most horizontal line and bottom-most edge of the triangle are parallel, so by the alternate interior angles theorem, $\angle x'=\angle x$ and $\angle y'=\angle y$. \ Therefore, the sum of the interior angles of the triangle ($x$, $y$, and $z$) is also $\pi $.

When $n=3$, we also see that $\pi ( n-2) =\pi $. Therefore $F( 3)$ is true.

\

For the induction step, we aim to prove that if we assume $F( k)$ is true for an arbitrary natural number $k\geq 3$, then $F( k+1)$ is also true.

Given a convex polygon with $k+1$ verticies, we connect the first and third verticies with a line segment. This forms a triangle between the first, second, and third verticies, and a convex polygon with $k$ verticies. Both of these are contained within the original polygon as per the given definition of a convex polygon (``for any two points in the interior, the line segment that joins them is contained within the polygon'').

As proven in our base case, the sum of the interior angles of the triangle is $\pi $. Given our assumption that $F( k)$ is true, the sum of the interior angles of the convex polygon with $k$ verticies will be $\pi (k-2)$. Because both of these polygons are contained within the original polygon, the sum of the interior angles of the entire convex polygon with $k+1$ verticies will be $\pi +\pi ( k-2) =\pi ( k-1)$.

Notice that this is equal to $\pi ( n-2)$ when $n=k+1$. Therefore, if we assume $F(k)$ is true for an arbitrary natural number $k\geq 3$, then $F(k+1)$ is also true.

\

Having proven that $F(3)$ is true, and that if we assume $F(k)$ is true for an arbitrary natural number $k\geq 3$ then $F(k+1)$ is also true, we can use the principle of induction to conclude that $F(n)$ is true for all natural numbers $n\geq 3$.\qed 

\
\hline
%%%%%%%%%%%%%%%%%%%%%%%%%%%%%%%%%%%%%%%%%%%%%%%%%%%%%%%%%%%%%%%%%%%%%%%%%%%%%%%%%
\section{}



\begin{problem*}
    The Computer Science department considers induction the most important part of Math 300 (even though it \textit{definitely} is not), and the reason why this is the case is that induction provides a way to prove that an algorithm is correct. Consider the following \texttt{Adamsort} algorithm, which takes an array of length $n$ as input and outputs an array which contains the same elements, but sorted. 
    
    \begin{center}
    	\begin{minipage}{0.7\linewidth} % Adjust the minipage width to accomodate for the length of algorithm lines
    		\begin{algorithm}[H]
                \textbf{Require:} $A$ an array of length $n \geq 0$
                
                \textbf{Ensure:} $A$ is sorted at the end. That is, $i < j$, then $A[i] \leq A[j]$.
    			
                \medskip
    			
                \textbf{for} $1 \leq i \leq \binom{n}{2}$ \textbf{do}
    			
                    \Indp 
                    \textbf{if} there exist $j < k$ such that $A[i] > A[k]$ \textbf{then} 
        			
        			    \Indp \Indp
            			Swap $A[j]$ and $A[k]$
            		
            		\Indm \Indm	
            		\textbf{else}
            		
            		    \Indp \Indp
            		    Break
            		    
            	    \Indm \Indm
            	    \textbf{end if}
            	
            	\Indm
            	\textbf{end for}
            	    
        	    \caption{\texttt{Adamsort}} % Algorithm name
    			\label{alg:Adamsort}   % optional label to refer to
    		\end{algorithm}
    	\end{minipage}
    \end{center}
    
    % \begin{algorithm}\caption{Adamsort}\begin{algorithmic}
    % \Require $A$ an array of length $n\ge 0$
    % \Ensure $A$ is sorted at the end. That is, $i<j$, then $A[i]\le A[j]$.
    % \For{$1\le i\le \binom{n}{2}$} 
    % \If {there exist $j<k$ such that $A[j]>A[k]$} 
    % \State Swap $A[j]$ and $A[k]$
    % \Else \State Break 
    % \EndIf
    % \EndFor
    
    % \end{algorithmic}\end{algorithm}
    
    Prove using induction that this algorithm converges to a sorted list.
\end{problem*}

Define the proposition
\begin{equation*}
F( n) :=\text{for an array} \ A\ \text{with} \ n\text{ elements, Adamsort converges to a sorted list} .
\end{equation*}
We aim to prove by induction that $F( n)$ is true for all integers $n$ greater than 0.

\

We consider the first 2 base cases. 

For $n=0$, $F( 0)$ is trivially true as there are no elements to sort.

For $n=1$, $F( 1)$ is also trivially true as an array with one element is already sorted.

\

For the induction step, we aim to prove that if we assume $F( k)$ is true for an arbitrary integer $k\geq 1$, then $F( k+1)$ is also true.

For an array with $k+1$ elements, our assumption that $F(k)$ is true means that Adamsort will be able to sort the first $k$ elements in $\binom{k}{2}$ iterations. In the best case scenario, the list is already sorted by this point, so the \textbf{else} branch is entered and terminates the algorithm. In the worst case scenario, the first element is greater than the last element ($A[ 1]  >A[ k+1]$), in which case the \textbf{if} branch is entered and iterates $k$ times before swapping the two elements, resulting in a sorted list. These $\binom{k}{2} + k$ total iterations are allowed within the maximum $\binom{k+1}{2}$ iterations of the \textbf{for} loop, as 
\begin{align*}
\binom{k+1}{2} & =\binom{( k+1) -1}{( 2)} +\binom{( k+1) -1}{( 2) -1} \quad \text{(Pascal's Identity, proven in 7a)}\\
 & =\binom{k}{2} +\binom{k}{1}\\
 & =\binom{k}{2} +\frac{k!}{1!( k-1) !}\\
 & =\binom{k}{2} +k.
\end{align*}
At this point, the \textbf{for} loop will have reached its max iterations and will terminate. In both the best and worst cases the algorithm converges to a sorted list, so $F( k+1)$ is true.

\

Having proven that $F(0)$ and $F( 1)$ are true, and that if we assume $F(k)$ is true for an arbitrary natural number $k\geq 1$ then $F(k+1)$ is also true, we can use the principle of induction to conclude that $F(n)$ is true for all integers $n\geq 2$.\qed 


\
\hline
%%%%%%%%%%%%%%%%%%%%%%%%%%%%%%%%%%%%%%%%%%%%%%%%%%%%%%%%%%%%%%%%%%%%%%%%%%%%%%%%%
\section{}

\begin{problem} a
\textbf{(Pascal's Identity).} The \textbf{binomial coefficient} is written \[\binom{n}{k}:=\frac{n!}{k!(n-k)!}\] (\verb|\binom{n}{k}|) and is the coefficient on $x^k$ of the expansion of the binomial $(1+x)^n$. It is also the number of ways to choose $k$ elements from a set with $n$ elements. Prove that \[\binom{n}{k}=\binom{n-1}{k}+\binom{n-1}{k-1}.\] [Hint: no need to use induction here, just verify using computation.]
\end{problem}

First, we note the factorial property that $x\cdot (x-1)!=x!$.

We proceed to compute:
\begin{align*}
    \binom{n-1}{k} +\binom{n-1}{k-1} & =\frac{( n-1) !}{k!( n-k-1) !} +\frac{( n-1) !}{( k-1) !( n-k) !}\\
     & =( n-1) !\left(\frac{1}{k!( n-k-1) !} +\frac{1}{( k-1) !( n-k) !}\right)\\
     & =( n-1) !\left(\frac{( n-k)}{k!( n-k) !} +\frac{k}{k!( n-k) !}\right)\\
     & =( n-1) !\left(\frac{n}{k!( n-k) !}\right)\\
 & =\frac{n!}{k!( n-k) !} .
\end{align*} \qed

\begin{problem} b
\textbf{(Leibniz Rule).} $\dagger$ Suppose $f$ and $g$ are differentiable functions. We know from high school calculus the product rule: \[(fg)'=f'g+fg'.\] Prove using induction the following Leibniz Rule: \[(fg)^{(n)}=\sum_{k=0}^n \binom{n}{k}f^{(n-k)}g^{(k)}.\] Here $f^{(n)}$ denotes the $n$th derivative of $f$, with the convention that $f^{(0)}=f$. [Hint: use part (a).]
\end{problem}

\begin{tikzpicture}
    \amongUsI[scale=0.1]{red}{gray}
\end{tikzpicture}

\
\hline
%%%%%%%%%%%%%%%%%%%%%%%%%%%%%%%%%%%%%%%%%%%%%%%%%%%%%%%%%%%%%%%%%%%%%%%%%%%%%%%%%
\end{document}